\chapter*{出版说明}
\addcontentsline{toc}{chapter}{出版说明}

为深入实施科教兴国战略、人才强国战略、创新驱动发展战略，统筹推进教育科技人才体制机制一体化改革，教育部于2023年4月19日正式启动基础学科系列本科教育教学改革试点工作（下称“101计划”）。物理学领域“101计划”工作组邀请国内物理学界教学经验丰富、学术造诣深厚的优秀教师和顶尖专家，及31所基础学科拔尖学生培养计划2.0基地建设高校，从物理学专业教育教学的基本规律和基础要素出发，共同探索建设一流核心课程、一流核心教材、一流核心教师团队和一流核心实践项目。这一系列举措有效地提高了我国物理学专业本科教学质量和水平，引领带动相关专业本科教育教学改革和人才培养质量提升。

通过基础要素建设的“小切口”，牵引教育教学模式的“大改革”，让人才培养模式从“知识为主”转向“能力为先”，是基础学科系列“101计划”的主要目标。物理学领域“101计划”工作组遴选了力学、热学、电磁学、光学、原子物理学、理论力学、电动力学、量子力学、统计力学、固体物理、数学物理方法、计算物理、实验物理、物理学前沿与科学思想选讲等14门基础和前沿兼备、深度和广度兼顾的一流核心课程，由课程负责人牵头，组织调研并借鉴国际一流大学的先进经验，主动适应学科发展趋势和新一轮科技革命对拔尖人才培养的要求，力求将“世界一流”“中国特色”“101风格”统一在配套的教材编写中。本教材系列在吸纳新知识、新理论、新技术、新方法、新进展的同时，注重推动弘扬科学家精神，推进教学理念更新和教学方法创新。

在教育部高等教育司的周密部署下，物理学领域“101计划”工作组下设的课程建设组、教材建设组，联合参与的教师、专家和高校，以及北京大学出版社、高等教育出版社、科学出版社等，经过反复研讨、协商，确定了系列教材详尽的出版规划和方案。为保障系列教材质量，工作组还专门邀请多位院士和资深专家对每种教材的编写方案进行评审，并对内容进行把关。

在此，物理学领域“101计划”工作组谨向教育部高等教育司的悉心指导、31所参与高校的大力支持、各参与出版社的专业保障表示衷心的感谢；向北京大学郝平书记、龚旗煌校长，以及北京大学教师教学发展中心、教务部等相关部门在物理学领域“101计划”酝酿、启动、建设过程中给予的亲切关怀、具体指导和帮助表示由衷的感谢；特别要向14位一流核心课程建设负责人及参与物理学领域“101计划”一流核心教材编写的各位教师的辛勤付出，致以诚挚的谢意和崇高的敬意！

基础学科系列“101计划”是我国本科教育教学改革的一项筑基性工程。改革，改到深处是课程，改到实处是教材。物理学领域“101计划”立足世界科技前沿和国家重大战略需求，以兼具传承经典和探索新知的课程、教材建设为引擎，着力推进卓越人才自主培养，激发学生的科学志趣和创新潜力，推动教师为学生成长成才提供学术引领、精神感召和人生指导。本教材系列的出版，是物理学领域“101计划”实施的标志性成果和重要里程碑，与其他基础要素建设相得益彰，将为我国物理学及相关专业全面深化本科教育教学改革、构建高质量人才培养体系提供有力支撑。

\begin{flushright}
物理学领域“101计划”工作组
\end{flushright}
\clearpage
